% ---------------------------------------------------------------------
%  第 9 章 Fourier 变换
% ---------------------------------------------------------------------
\chapter{Fourier 变换}
\label{ch:fourier}

库中的 Fourier 变换有三类用途：把虚时间数据变换到 Matsubara 频率、
把稳态数据在时间与频率之间往返、以及在自适应求积中对一般函数做
Fourier 积分。三者都建立在一个自写的任意长度 FFT 之上~\cite{cooley1965,bluestein1970}。本章先说明
FFT 的算法与缓存，再说明三种变换的离散化与修正。

\section{FFT 算法}
\label{sec:fft}

\cref{fig:mixed-radix}以一个具体长度为例，展示混合基数变换的分解顺序。

离散 Fourier 变换定义为
\begin{equation}
    X_k = \sum^{N-1}_{n=0} x_n\,\ee^{2\pi\ii s nk/N},
    \qquad s=-1\ (\text{正向}),\quad s=+1\ (\text{反向}),
    \label{eq:dft}
\end{equation}
两个方向都不做归一化，因此正向接反向会把数据乘以 $N$。

\paragraph{混合基数分解。}
若 $N$ 的所有素因子都不超过 $13$，采用自排序的 Stockham
Cooley--Tukey 算法。设当前长度为 $n$、基为 $p$、$m=n/p$，
以及此前各基的乘积 $s'$，一级蝶形为
\begin{equation}
    y_{q+s'(pk+u)} =
    W^{uk}_{n}\sum^{p-1}_{t=0} W^{ut}_{p}\,
      x_{q+s'(k+tm)},
    \qquad W_n = \ee^{2\pi\ii s/n}.
    \label{eq:stockham}
\end{equation}
基的优先顺序是 $4,2,3,5,7,11,13$：以 $4$ 为主可以减少级数与
临时存储，$2$ 用来吸收奇偶性，其余素因子按从大到小排出。

\begin{figure}[htbp]
    \centering
    \begin{tikzpicture}[
        node distance=6mm,
        box/.style={draw=kbblue!60, fill=kbblue!5, rounded corners=2pt,
                    font=\small, inner sep=4pt},
        ar/.style={-{Stealth[length=5pt]}, thick, kbblue!60}
    ]
        \node[box] (n) {$N=2^3\cdot 3\cdot 5$};
        \node[box, below=of n] (s1) {基 $4$：$N/4$ 组};
        \node[box, below=of s1] (s2) {基 $2$：$N/2$ 组};
        \node[box, below=of s2] (s3) {基 $3$：$N/3$ 组};
        \node[box, below=of s3] (s4) {基 $5$：$N/5$ 组};
        \node[box, below=of s4, fill=kbgreen!6, draw=kbgreen!60]
            (out) {输出，自然顺序};

        \draw[ar] (n) -- (s1);
        \draw[ar] (s1) -- (s2);
        \draw[ar] (s2) -- (s3);
        \draw[ar] (s3) -- (s4);
        \draw[ar] (s4) -- (out);
    \end{tikzpicture}
    \caption[混合基数分解]{混合基数分解。每一级把长度缩小一个基因子；Stockham
             算法把结果写到另一组缓冲区，因此最后自然得到正常顺序，
             不需要位反转。}
    \label{fig:mixed-radix}
\end{figure}

\section{Bluestein 回退}
\label{sec:bluestein}

若 $N$ 含大于 $13$ 的素因子，混合基数分解不再适用，改用 Bluestein
变换。利用
\begin{equation}
    nk = \frac{n^2+k^2-(k-n)^2}{2},
\end{equation}
把 \cref{eq:dft} 写成卷积：
\begin{equation}
    X_k = c_k \sum^{N-1}_{n=0} (x_n c_n)\,b_{k-n},
    \qquad
    c_n = \ee^{\ii\pi s n^2/N},\quad
    b_m = \ee^{-\ii\pi s m^2/N}.
    \label{eq:bluestein}
\end{equation}
卷积用长度为 $M=2^{\lceil\log_2(2N-1)\rceil}$ 的循环卷积实现，
三次 FFT 完成。

\section{plan 缓存}
\label{sec:plan-cache}

FFT 的旋转因子表与基分解只依赖于长度与方向，因此按
$(N,\text{方向})$ 缓存：
\begin{lstlisting}[numbers=none]
pub fn plan_fft(n: usize, direction: FftDirection) -> Arc<FftPlan>
\end{lstlisting}
返回的 \code{Arc<FftPlan>} 可以跨调用共享；内部用
\code{OnceLock<Mutex<HashMap<...>>>} 保存。缓存在多次变换之间
复用旋转因子，避免重复计算三角函数。\code{clear_fft_plan_cache}
用于在内存敏感的场景中释放缓存。

\section{批量变换}
\label{sec:strided-fft}

稳态数据与 Matsubara 数据的布局都是“轨道指标在前，模式索引在后”，
因此需要沿着一个带步长的轴做多次变换。\code{fft_strided} 的语义是：
输入是一段连续数据，变换沿着指标集合
\begin{equation}
    \{\,i_0 + m\cdot\text{stride} \mid m=0,\dots,N-1\,\}
\end{equation}
进行，其余指标视为批量维度。它一次性完成所有批次的变换，避免逐个
序列调用带来的函数开销。\code{fft_strided_in_place} 是它的就地版本。

\section{离散修正}
\label{sec:dft-correction}

\cref{tab:corr}比较不同端点修正的误差阶与网格要求。

直接对等距采样做离散 Fourier 变换（Discrete Fourier Transform，DFT），其积分近似误差通常为
$O(\delta)$。为提高精度，库采用端点修正（end correction）：
把积分写成
\begin{equation}
    I(\omega) \approx \delta\Big\{
    W(\theta)\sum^{N}_{r=0}\ee^{\ii\omega t_r}f_r
    +\sum^{3}_{j=0}\big[\alpha_j(\theta)f_j
      +\ee^{\ii\omega N\delta}\alpha^{*}_j(\theta)f_{N-j}\big]\Big\},
    \label{eq:endcor}
\end{equation}
其中 $\theta=\omega\delta$，$W$ 是修正因子，$\alpha_j$ 是边界项。
线性修正只保留 $j=0,1$ 两项，三次修正保留四项。

\begin{table}[htbp]
    \centering
    \caption[两种修正的适用场合]{两种修正的适用场合}
    \label{tab:corr}
    \small
    \begin{tabularx}{\textwidth}{@{}llX@{}}
        \toprule
        方法 & 函数 & 特点 \\
        \midrule
        线性 & \code{dftcorr_linear} &
            $W(\theta)$ 与两个边界项；小 $\theta$ 用级数展开 \\
        三次 & \code{dftcorr_cubic} &
            四个边界项；误差为 $O(\delta^4)$；小 $\theta$ 用级数展开 \\
        面板 & \code{dft_panel} &
            在一段区间上对三次插值精确积分，供自适应算法使用 \\
        \bottomrule
    \end{tabularx}
\end{table}

当 $|\theta|<0.05$ 时，直接用三角函数的表达式会出现相消误差，
因此代码改用 Taylor 级数；两个分支在 $\theta=0.05$ 附近给出的值
一致到机器精度。

\section{Matsubara 变换}
\label{sec:matsubara-fourier}

Matsubara 频率为
\begin{equation}
    \omega_m = \frac{\pi}{\bet}(2m+\delta_\xi),
    \qquad
    \delta_\xi =
    \begin{cases}1, & \text{费米型},\\ 0, & \text{玻色型}.\end{cases}
    \label{eq:matsubara-freq}
\end{equation}
\code{matsubara_dft} 计算未乘 $\Delta\tau$ 的原始变换：
\begin{equation}
    \tilde{A}_m = \sum^{N}_{r=0}
    \ee^{\ii\pi(2m+\delta_\xi)r/N}A(\tau_r),
    \qquad m=0,\dots,N,
\end{equation}
其中 $N=n_\tau$。\code{matsubara_ft} 在它之上加入端点修正与
$\Delta\tau$ 因子。

\paragraph{一阶尾部修正。}
虚时间 Green 函数在 $\tau=0,\bet$ 处一般有跳变。若直接变换，
跳变会以 $1/\omega$ 的慢衰减污染高频。库先减去一阶尾部
\begin{equation}
    c_0(\tau) =
    \begin{cases}
        -\dfrac{a_0}{2}, & \text{费米型},\\[4pt]
        a_0\Big(\dfrac{\tau}{\bet}-\dfrac12\Big), & \text{玻色型},
    \end{cases}
    \label{eq:tail}
\end{equation}
其中 $a_0$ 是端点跳变系数。修正后的函数连续，变换结果衰减更快，
逆变换后再把 $c_0$ 加回。

\paragraph{频域预解算。}
\code{MatsubaraSeed} 把“构造变换表 + 在每个频点上解一个
$d\times d$ 线性系统”封装成一个可复用的对象：
\begin{lstlisting}[numbers=none]
let mut seed = MatsubaraSeed::new(&f.mat, beta, sign, pcf, order);
seed.solve(&mut x, &b);
\end{lstlisting}
参数 \code{pcf} 是过采样窗口的份数，\code{order} 是修正阶数，
取 $0$、$1$ 或 $3$，分别对应无修正、线性修正与三次修正。由于核在所有迭代中不变，种子只需构造一次；
时间推进与迭代求解器都利用了这一点。

\section{自适应采样}
\label{sec:adaptive}

对没有解析结构的函数，库提供自适应 Fourier 积分。核心是一个
“面板”结构：区间 $[a,b]$ 被切成若干段，每段存 $n_n+1$ 个等距
采样值；对每个面板用三次修正的 DFT 积分，再用“隔点比较”估计误差。
若误差超过容差，就把该面板一分为二，递归细化。

\begin{lstlisting}[numbers=none]
pub fn adft_func(f: &dyn Fn(f64) -> Complex, a: f64, b: f64,
                 omega: f64, limit: i64) -> (Complex, Complex)
\end{lstlisting}
\code{limit} 是面板数的上限。函数返回积分值与误差估计。
\code{dft_cplx} 要求偶数个、且大于 $16$ 的采样点；这是因为三次
修正需要成对的端点，而级数展开在点数过少时不收敛。

\section{稳态变换}
\label{sec:steady-transform}

稳态数据在两个域中各有一份系数。频域变换是
\code{integral_transform_to_freq}，时域变换是
\code{integral_transform_to_time}。两者都支持：
\begin{itemize}
    \item \code{Trapez}：梯形法则，等价于对周期数据做普通 DFT；
    \item \code{Cubic}：三次修正，精度更高；
    \item \code{Adaptive}：只对平衡态构造有意义，积分变换本身不支持。
\end{itemize}

\paragraph{对称性。}
推迟分量在 $t<0$ 时为零，超前分量满足
$\Gadv(-t)=\big[\Gret(t)\big]^{\dagger}$。小于分量满足
$\Gless(-t)=-\big[\Gless(t)\big]^{\dagger}$，因此可由非负时间样点重构负时间半轴。
时间网格的 Nyquist 槽在物理积分变换中按实现约定置零。

\section{轮廓与稳态互转}
\label{sec:convert}

\paragraph{重采样。}
\code{downsample}、\code{upsample} 与 \code{resample} 在稳态网格上
改变分辨率。下采样每隔 $f$ 个点抽取一个时间样点，并保留较低的频率模式。
上采样先用梯形法则由时间数据重新生成频谱，再对新添的频率模式补零并逆变换。
两者改变时间步长，但保持频率间距不变。

\paragraph{由轮廓提取稳态。}
\code{contour_to_steady} 固定较晚的实时间 $t_n$，读取推迟块
$G^R(t_n,t_{n-j})$ 和小于块 $G^<(t_{n-j},t_n)$。取样范围为
$0\le j\le\min(n,N_{\mathrm{ft}}/2-1)$，其余未覆盖的时间模式清零。
相反符号的相对时间小于块由伴随关系补齐。该操作只更新时域数组。
\code{steady_to_contour} 按时间平移不变性回填实时间的推迟和小于分量，
超出稳态时间窗的块置零。虚时和左混合分量保留原值。
